\section{USAMO 2008}
        \begin{enumerate}
        	\item (\textit{USAMO 2008}) (\textit{Titu Andreescu}) Prove that for each positive integer $n$, there are pairwise relatively prime integers $k_0,k_1\ldots,k_n$, all strictly greater than 1, such that $k_0k_1\cdots k_n-1$ is the product of two consecutive integers.


 

	\item (\textit{USAMO 2008}) (\textit{Zuming Feng}) Let $ABC$ be an acute, \href{https://artofproblemsolving.com/wiki/index.php?title=Scalene}{scalene} triangle, and let $M$, $N$, and $P$ be the midpoints of $\overline{BC}$, $\overline{CA}$, and $\overline{AB}$, respectively. Let the \href{https://artofproblemsolving.com/wiki/index.php?title=Perpendicular}{perpendicular} \href{https://artofproblemsolving.com/wiki/index.php?title=Bisect}{bisectors} of $\overline{AB}$ and $\overline{AC}$ intersect ray $AM$ in points $D$ and $E$ respectively, and let lines $BD$ and $CE$ intersect in point $F$, inside of triangle $ABC$. Prove that points $A$, $N$, $F$, and $P$ all lie on one circle.


 

	\item (\textit{USAMO 2008}) (\textit{Gabriel Carroll}) Let $n$ be a positive integer. Denote by $S_n$ the set of points $(x, y)$ with integer coordinates such that

 \[\left|x\right|+\left|y+\frac{1}{2}\right|<n\]
A path is a sequence of distinct points $(x_1,y_1),(x_2,y_2),\ldots,(x_\ell,y_\ell)$ in $S_n$ such that, for $i=2,\ldots,\ell$, the distance between $(x_i,y_i)$ and $(x_{i-1},y_{i-1})$ is $1$ (in other words, the points $(x_i,y_i)$ and $(x_{i-1},y_{i-1})$ are neighbors in the lattice of points with integer coordinates). Prove that the points in $S_n$ cannot be partitioned into fewer than $n$ paths (a partition of $S_n$ into $m$ paths is a set $\mathcal{P}$ of $m$ nonempty paths such that each point in $S_n$ appears in exactly one of the $m$ paths in $\mathcal{P}$).


 

	\item (\textit{USAMO 2008}) (\textit{Gregory Galparin}) Let $\mathcal{P}$ be a \href{https://artofproblemsolving.com/wiki/index.php?title=Convex\_polygon}{convex polygon} with $n$ sides, $n\ge3$. Any set of $n-3$ diagonals of $\mathcal{P}$ that do not intersect in the interior of the polygon determine a \textit{triangulation} of $\mathcal{P}$ into $n - 2$ triangles. If $\mathcal{P}$ is regular and there is a triangulation of $\mathcal{P}$ consisting of only isosceles triangles, find all the possible values of $n$.


 

	\item (\textit{USAMO 2008}) (\textit{Kiran Kedlaya}) Three nonnegative real numbers $r_1$, $r_2$, $r_3$ are written on a blackboard. These numbers have the property that there exist integers $a_1$, $a_2$, $a_3$, not all zero, satisfying $a_1r_1+a_2r_2+a_3r_3=0$. We are permitted to perform the following operation: find two numbers $x$, $y$ on the blackboard with $x\le y$, then erase $y$ and write $y-x$ in its place. Prove that after a finite number of such operations, we can end up with at least one $0$ on the blackboard.


 

	\item (\textit{USAMO 2008}) (<i>\href{https://artofproblemsolving.com/wiki/index.php?title=Sam\_Vandervelde}{Sam Vandervelde}</i>) At a certain mathematical conference, every pair of mathematicians are either friends or strangers. At mealtime, every participant eats in one of two large dining rooms. Each mathematician insists upon eating in a room which contains an even number of his or her friends. Prove that the number of ways that the mathematicians may be split between the two rooms is a power of two (i.e., is of the form $2^k$ for some positive integer $k$).


 

\end{enumerate}